\documentclass[twocolumn,10pt]{article}

\usepackage{fontspec}
\usepackage{polyglossia}
\setdefaultlanguage{russian}
\setotherlanguage{english}
\setmainfont{Times New Roman}
\usepackage{amsmath,amssymb}
\usepackage{booktabs}
\usepackage{graphicx}
\usepackage[margin=2cm]{geometry}
\usepackage{xcolor}
\usepackage{hyperref}

\hypersetup{
  colorlinks=true,
  linkcolor=blue!60!black,
  citecolor=blue!60!black,
  urlcolor=blue!60!black,
  pdfauthor={David Alfyorov},
  pdftitle={Absence of Starobinsky inflation in the spectral action
    at conformal coupling}
}

\newcommand{\dd}{\mathrm{d}}
\newcommand{\erfi}{\operatorname{erfi}}
\newcommand{\order}{\mathcal{O}}

\begin{document}

\twocolumn[
\begin{center}
{\Large\bfseries Отсутствие инфляции Старобинского\\
в спектральном действии при конформной связи}

\bigskip
{\large Д.\,Алфёров}

\smallskip
{\small\textit{davidich.alfyorov@gmail.com}}

\bigskip
\begin{minipage}{0.85\textwidth}
\small\textbf{Аннотация.}
Стандартная спектральная тройка Шамседдина--Конна
предсказывает конформную связь $\xi = 1/6$ для поля Хиггса.
Показано, что это предсказание устраняет член $R^2$ из
однопетлевого эффективного действия: коэффициент
$\alpha_R = 2(\xi - 1/6)^2$ обращается в нуль при $\xi = 1/6$.
Скалярон отсутствует в гравитационном спектре, и инфляция
Старобинского исключена в рамках стандартного спектрального
действия. Рассмотрены восемь механизмов восстановления скалярона;
все терпят неудачу. Теория предсказывает ровно две поляризации
гравитационных волн. Обнаружение скалярной поляризации
фальсифицировало бы стандартную спектральную тройку.
\end{minipage}

\bigskip
\end{center}
]

\section*{1.}
Спектральное действие
$S = \mathrm{Tr}\,f(D^2/\Lambda^2)$~\cite{CC:1996,CC:1997}
выводит бозонное действие гравитации и Стандартной Модели из одного
операторного следа. Разложение по ядру теплопроводности
даёт~\cite{Vassilevich:2003} действие
Эйнштейна--Гильберта, космологическую постоянную и квадратичные
по кривизне поправки с коэффициентами, определяемыми коэффициентами
Сили--Де~Витта~$a_{2k}$.

Квадратично-кривизнный сектор однопетлевого эффективного действия
имеет вид~\cite{BV:1987,BV:1990,CZ:2012}
\begin{equation}\label{eq:Gamma}
  \Gamma^{(1)} \supset \int \dd^4x\,\sqrt{g}\,
  \bigl[\alpha_C\, C_{\mu\nu\rho\sigma}C^{\mu\nu\rho\sigma}
  + \alpha_R(\xi)\, R^2\bigr]\,\ln\frac{\Lambda^2}{\mu^2}\,,
\end{equation}
где $\alpha_C = 13/120$ зависит только от содержания Стандартной
Модели ($N_s = 4$, $N_D = 45/2$, $N_v = 12$) и не зависит ни
от~$\xi$, ни от обрезающей функции~$f$~\cite{Alfyorov:2026paper1},
а коэффициент при~$R^2$ равен
\begin{equation}\label{eq:alphaR}
  \alpha_R(\xi) = 2\Bigl(\xi - \frac{1}{6}\Bigr)^{\!2}.
\end{equation}

\section*{2.}
В спектральной тройке Шамседдина--Конна хиггсовский дублет
возникает из конечной части оператора Дирака, и связь с кривизной
$\xi R|H|^2$ не является свободным параметром. Коэффициент~$a_4$
содержит общий юкава-зависимый множитель перед связью с кривизной
и кинетическим членом~$|\nabla H|^2$. После канонической
нормировки получается~\cite{CCM:2006,CC:2010,CCS:2013,vS:2015,DLM:2014}
\begin{equation}\label{eq:xi}
  \xi = \frac{1}{6}\qquad\text{(конформная связь)}.
\end{equation}
Этот результат подтверждён пятью независимыми выводами в рамках
некоммутативной геометрии, включая расширение
Пати--Салама~\cite{CCS:2013ps}, и устойчив при добавлении
скаляров из конечного оператора Дирака~\cite{BFS:2010,vdD:2018}.

Однопетлевая бета-функция
$\beta_\xi \propto (\xi - 1/6)$~\cite{ParkerToms:2009}
обращается в нуль при $\xi = 1/6$. Полное однопетлевое уравнение
ренормгруппы~\cite{ParkerToms:2009,Buchbinder:1992}
\begin{equation}\label{eq:betaxi}
  \beta_\xi = \frac{1}{(4\pi)^2}\Bigl(\xi - \frac{1}{6}\Bigr)
  \bigl(6\lambda + 6y_t^2 - \tfrac{3}{2}g'^{\,2}
  - \tfrac{9}{2}g^2\bigr)
\end{equation}
имеет множитель $(\xi - 1/6)$, обеспечивающий, что $\xi = 1/6$
является неподвижной точкой при любых значениях констант связи СМ.

\section*{3.}
Следствия непосредственны. При $\xi = 1/6$:
\begin{equation}\label{eq:alphaR-zero}
  \alpha_R\bigl(\xi = \tfrac{1}{6}\bigr) = 0\,.
\end{equation}
Член $R^2$ отсутствует в однопетлевом эффективном действии.
Скалярона нет. Инфляция Старобинского~\cite{Starobinsky:1980},
требующая распространяющейся скалярной степени свободы из
члена~$R^2$ с массой $M_{\mathrm{inf}} \approx 1.28 \times
10^{-5}\,M_P$~\cite{Planck:2018}, исключена в стандартном
спектральном действии.

Знаменатель скалярного пропагатора
\begin{equation}\label{eq:Pis}
  \Pi_s(z, \xi) = 1 + 6\Bigl(\xi - \frac{1}{6}\Bigr)^{\!2}\,z\,
  \hat{F}_2(z, \xi)
\end{equation}
тождественно равен единице при $\xi = 1/6$, для любого
импульса~$z = k^2/\Lambda^2$ и любой обрезающей функции~$f$.
Скалярная гравитационная мода не распространяется.
Гравитационный спектр содержит только безмассовый спин-2
гравитон и массивный спин-2 факеон с массой
$m_2 = \sqrt{z_0}\,\Lambda$~\cite{Alfyorov:2026paper1}.
Скорость гравитационных волн $c_T = c$
точно~\cite{Alfyorov:2026paper4}, что совместимо с ограничением
GW170817: $|c_T - c|/c < 10^{-15}$~\cite{GW170817:2017}.

\section*{4.}
Можно задаться вопросом, возможно ли восстановить скалярон
в рамках спектрального действия. В таблице рассмотрены
восемь механизмов; все терпят неудачу.

\begin{table}[ht]
\centering
\caption{Механизмы восстановления инфляции Старобинского
  в спектральном действии и причина неудачи каждого.}
\label{tab:escape}
\begin{tabular}{p{3.5cm}p{4.8cm}}
\toprule
Механизм & Препятствие \\
\midrule
Субпланковский~$\Lambda$ &
  Конфликт с GUT-интерпретацией~$\Lambda$ \\
Большой $\xi \sim 2 \times 10^4$ &
  Нарушает геометрические ГУ спектральной тройки \\
$\sigma$-синглет НКГ~\cite{BFS:2010} &
  $\xi_\sigma = 1/6$ (конформная) \\
Пати--Салам~\cite{CCS:2013ps} &
  Все скаляры конформны \\
Grand Symmetry~\cite{CCS:2013gs} &
  $\sigma$ для массы Хиггса, не для скалярона \\
$N_s' \sim 4{\times}10^{10}$ &
  Требует содержания, несовместимого с НКГ \\
Модификация~$f$ &
  $\alpha_R$ из~$a_4$, не зависит от~$f$ \\
Двухпетлевые поправки &
  Вычисление отсутствует \\
\bottomrule
\end{tabular}
\end{table}

Единственные известные пути --- выход за рамки стандартного
спектрального действия: дзета-спектральное
действие~\cite{KLV:2015} или дилатонизированное
действие~\cite{CC:2006dilaton}. Они модифицируют операторный
след и не могут быть получены из стандартного
$\mathrm{Tr}\,f(D^2/\Lambda^2)$.

\section*{5.}
Результат даёт чёткий критерий фальсификации. Обнаружение
скалярной поляризации гравитационных волн (например, продольной
``дышащей'' моды в данных LIGO/Virgo/KAGRA или LISA) означало бы
$\xi \neq 1/6$ и фальсифицировало бы стандартную спектральную
тройку Шамседдина--Конна. Текущие ограничения из хронометража
двойных пульсаров~\cite{Stairs:2003} и события
GW170814~\cite{GW170814:2017} совместимы с отсутствием скалярного
содержания.

При $\xi = 1/6$ отношение $c_1/c_2 = -1/3$ в базисе Стелле
$\{R^2, R_{\mu\nu}^2\}$ заморожено вдоль однопетлевой траектории
ренормгруппы:
$\beta_{c_1}/\beta_{c_2} = -1/3$
тождественно~\cite{Alfyorov:2026paper1}. Это отношение
отличает SCT от асимптотической безопасности
($c_1/c_2 \approx -1.09$, зависит от
усечения~\cite{BMS:2009}) и от пертурбативной ветви асимптотической
свободы ($c_1/c_2 \approx -0.326$~\cite{Shapiro:2004}).

\section*{6.}
Итого: стандартное спектральное действие предсказывает
конформную связь $\xi = 1/6$, являющуюся однопетлевой
неподвижной точкой ренормгруппы. При $\xi = 1/6$ коэффициент
$\alpha_R$ обращается в нуль, скалярная гравитационная мода
отщепляется, и инфляция Старобинского исключена. Восемь
механизмов восстановления скалярона терпят неудачу в рамках
стандартной спектральной тройки. Предсказание ровно двух
поляризаций гравитационных волн проверяемо текущими и
планируемыми детекторами.

\bigskip

\noindent\textbf{ФИНАНСИРОВАНИЕ РАБОТЫ}

Данная работа выполнена без внешнего финансирования.
Никаких дополнительных грантов на проведение
или руководство данным конкретным исследованием получено не было.

\bigskip

\noindent\textbf{КОНФЛИКТ ИНТЕРЕСОВ}

Автор заявляет об отсутствии конфликта интересов.

\begin{thebibliography}{99}

\bibitem{CC:1996}
A.\,H.~Chamseddine, A.~Connes,
``The spectral action principle'',
Commun.\ Math.\ Phys.\ \textbf{186}, 731 (1997);
hep-th/9606001.

\bibitem{CC:1997}
A.\,H.~Chamseddine, A.~Connes,
``Inner fluctuations of the spectral action'',
J.\ Math.\ Phys.\ \textbf{47}, 063504 (2006);
hep-th/0512169.

\bibitem{Vassilevich:2003}
D.\,V.~Vassilevich,
``Heat kernel expansion: user's manual'',
Phys.\ Rept.\ \textbf{388}, 279 (2003);
hep-th/0306138.

\bibitem{BV:1987}
A.\,O.~Barvinsky, G.\,A.~Vilkovisky,
``The generalized Schwinger--DeWitt technique in gauge theories
and quantum gravity'',
Phys.\ Rept.\ \textbf{119}, 1 (1985).

\bibitem{BV:1990}
A.\,O.~Barvinsky, G.\,A.~Vilkovisky,
``Covariant perturbation theory.~II. Second order in the
curvature. General algorithms'',
Nucl.\ Phys.\ B \textbf{333}, 471 (1990).

\bibitem{CZ:2012}
A.~Codello, O.~Zanusso,
``On the non-local heat kernel expansion'',
J.\ Math.\ Phys.\ \textbf{54}, 013513 (2013);
arXiv:1203.2034.

\bibitem{Alfyorov:2026paper1}
D.~Alfyorov,
``Nonlocal one-loop form factors of the spectral action
with Standard Model content'' (2026);
doi:10.22541/au.177507193.33027854/v1.

\bibitem{CCM:2006}
A.\,H.~Chamseddine, A.~Connes, M.~Marcolli,
``Gravity and the standard model with neutrino mixing'',
Adv.\ Theor.\ Math.\ Phys.\ \textbf{11}, 991 (2007);
hep-th/0610241.

\bibitem{CC:2010}
A.\,H.~Chamseddine, A.~Connes,
``Noncommutative geometry as a framework for unification of
all fundamental interactions including gravity. Part~I'',
Fortschr.\ Phys.\ \textbf{58}, 553 (2010);
arXiv:0901.0577.

\bibitem{CCS:2013}
A.\,H.~Chamseddine, A.~Connes, W.\,D.~van Suijlekom,
``Beyond the spectral standard model: emergence of
Pati--Salam unification'',
JHEP \textbf{11}, 132 (2013);
arXiv:1304.8050.

\bibitem{vS:2015}
W.\,D.~van Suijlekom,
\textit{Noncommutative Geometry and Particle Physics}
(Springer, 2015).

\bibitem{DLM:2014}
K.~van den Dungen, W.\,D.~van Suijlekom,
``Particle physics from almost-commutative spacetimes'',
Rev.\ Math.\ Phys.\ \textbf{24}, 1230004 (2012);
arXiv:1204.0328.

\bibitem{CCS:2013ps}
A.\,H.~Chamseddine, A.~Connes, W.\,D.~van Suijlekom,
``Inner fluctuations in noncommutative geometry without the
first order condition'',
J.\ Geom.\ Phys.\ \textbf{73}, 222 (2013);
arXiv:1304.7583.

\bibitem{BFS:2010}
T.\,D.~Brennan, S.\,E.~Flood, D.\,G.~Robbins,
``The $\sigma$-field in the Chamseddine--Connes spectral action'',
Eur.\ Phys.\ J.\ C \textbf{74}, 2862 (2014);
arXiv:1401.1006.

\bibitem{vdD:2018}
K.~van den Dungen,
``Families of spectral triples and foliations of space(time)'',
J.\ Noncommut.\ Geom.\ \textbf{12}, 1155 (2018);
arXiv:1711.07299.

\bibitem{CCS:2013gs}
A.\,H.~Chamseddine, A.~Connes, W.\,D.~van Suijlekom,
``Grand unification in the spectral Pati--Salam model'',
Fortschr.\ Phys.\ \textbf{62}, 797 (2014);
arXiv:1409.7574.

\bibitem{ParkerToms:2009}
L.~Parker, D.~Toms,
\textit{Quantum Field Theory in Curved Spacetime}
(Cambridge Univ.\ Press, 2009).

\bibitem{Buchbinder:1992}
И.\,Л.~Бухбиндер, С.\,Д.~Одинцов, И.\,Л.~Шапиро,
\textit{Эффективное действие в квантовой гравитации}
(Наука, М., 1992)
[I.\,L.~Buchbinder, S.\,D.~Odintsov, I.\,L.~Shapiro,
\textit{Effective Action in Quantum Gravity}
(IOP Publishing, Bristol, 1992)].

\bibitem{Starobinsky:1980}
А.\,А.~Старобинский,
``Новый тип изотропных космологических моделей без сингулярности'',
Письма в Астрон.\ ж.\ \textbf{6}, 681 (1980)
[A.\,A.~Starobinsky,
``A new type of isotropic cosmological models without singularity'',
Phys.\ Lett.\ B \textbf{91}, 99 (1980)].

\bibitem{Planck:2018}
Y.~Akrami и др.\ (Planck Collaboration),
``Planck 2018 results.~X. Constraints on inflation'',
Astron.\ Astrophys.\ \textbf{641}, A10 (2020);
arXiv:1807.06211.

\bibitem{Alfyorov:2026paper4}
D.~Alfyorov,
``Nonlinear field equations and FLRW cosmology
from the spectral action'' (2026);
doi:10.5281/zenodo.19098027.

\bibitem{GW170817:2017}
B.\,P.~Abbott и др.\ (LIGO/Virgo),
``Gravitational waves and gamma-rays from a binary neutron star
merger: GW170817 and GRB 170817A'',
Astrophys.\ J.\ Lett.\ \textbf{848}, L13 (2017);
arXiv:1710.05834.

\bibitem{KLV:2015}
B.~Iochum, C.~Levy, D.~Vassilevich,
``Spectral action beyond the weak-field approximation'',
Commun.\ Math.\ Phys.\ \textbf{340}, 1163 (2015);
arXiv:1503.09054.

\bibitem{CC:2006dilaton}
A.\,H.~Chamseddine, A.~Connes,
``Scale invariance in the spectral action'',
J.\ Math.\ Phys.\ \textbf{47}, 063504 (2006);
arXiv:0706.3690.

\bibitem{Stairs:2003}
I.\,H.~Stairs,
``Testing general relativity with pulsar timing'',
Living Rev.\ Relativ.\ \textbf{6}, 5 (2003).

\bibitem{GW170814:2017}
B.\,P.~Abbott и др.\ (LIGO/Virgo),
``GW170814: a three-detector observation of gravitational waves
from a binary-black-hole coalescence'',
Phys.\ Rev.\ Lett.\ \textbf{119}, 141101 (2017);
arXiv:1709.09660.

\bibitem{BMS:2009}
D.~Benedetti, P.\,F.~Machado, F.~Saueressig,
``Asymptotic safety in higher-derivative gravity'',
Nucl.\ Phys.\ B \textbf{824}, 168 (2010);
arXiv:0902.4630.

\bibitem{Shapiro:2004}
I.\,L.~Shapiro,
``Effective action of vacuum: semiclassical approach'',
Class.\ Quant.\ Grav.\ \textbf{25}, 103001 (2008);
arXiv:0801.0216.

\end{thebibliography}

\end{document}
